\chapter{Linear Algebra: Notation, Vectors, Matrices and Operations}
\label{chap:linalg}
%=============================================================================

\section{Notations and conventions}
\label{sec:notation}

Throughout these notes vectors, matrices and higher-order tensors are
always set in boldface.  Vectors are denoted by lower-case letters
($\bm{x}$, $\bm{v}$, \ldots) and matrices and tensors by upper-case
letters ($\bm{A}$, $\bm{U}$, \ldots).

Unless stated otherwise the elements $v_i$ of a vector $\bm{v}$ are
assumed real.  A real vector of length $n$ satisfies $\bm{x}\in
\mathbb{R}^{n}$; for complex vectors we write $\bm{x}\in
\mathbb{C}^{n}$.  A real $n\times n$ matrix satisfies $\bm{A}\in
\mathbb{R}^{n\times n}$.  Index counting starts at zero: matrix
elements run from $a_{00}$ to $a_{n-1,n-1}$.

We use the standard shorthand symbols: $\forall$ (for all),
$\Rightarrow$ (implies), $\equiv$ (equivalent / defined as), with
$\mathbb{R}$, $\mathbb{I}$, $\mathbb{C}$ for the real, integer and
complex number fields.

%----------------------------------------------------------------------
\section{Vectors}
\label{sec:vectors}

A column vector of length $n$ and its transpose are
\[
  \bm{x} = \begin{bmatrix} x_0 \\ x_1 \\ \vdots \\ x_{n-1} \end{bmatrix},
  \qquad
  \bm{x}^{T} =
    \begin{bmatrix} x_0 & x_1 & \cdots & x_{n-1} \end{bmatrix}.
\]
For complex vectors the \emph{Hermitian conjugate} (adjoint) is
\[
  \bm{x}^{\dagger} =
    \begin{bmatrix} x_0^* & x_1^* & \cdots & x_{n-1}^* \end{bmatrix}.
\]
The \emph{inner (dot) product} is the scalar
\[
  \bm{x}^{T}\bm{x} = \sum_{i=0}^{n-1} x_i x_i
    = x_0^2 + x_1^2 + \cdots + x_{n-1}^2,
\]
and for complex vectors
$\bm{x}^{\dagger}\bm{x} = \sum_{i=0}^{n-1} x_i^* x_i = \|{\bm x}\|^2$.

The \emph{outer product} of two real vectors $\bm{y}$ (length $n$)
and $\bm{z}$ (length $m$) is the $n\times m$ matrix
\[
  \bm{y}\bm{z}^{T} \;\Longrightarrow\; (\bm{y}\bm{z}^T)_{ij} = y_i z_j.
\]
For complex vectors one writes $\bm{y}\bm{z}^{\dagger}$ with
$(\bm{y}\bm{z}^{\dagger})_{ij}=y_i z_j^*$.

Important vector operations:
\begin{align*}
  \text{Addition / subtraction:}\quad
      &\bm{x} = \bm{y}\pm\bm{z} \;\Rightarrow\; x_i = y_i\pm z_i.\\
  \text{Scalar multiplication:}\quad
      &\bm{x} = \gamma\bm{y} \;\Rightarrow\; x_i = \gamma y_i.\\
  \text{Hadamard (element-wise) product:}\quad
      &\bm{x} = \bm{y}\circ\bm{z} \;\Rightarrow\; x_i = y_i z_i.
\end{align*}

%----------------------------------------------------------------------
\section{Matrices: definitions and algebraic operations}
\label{sec:matrices}

A general $n\times n$ matrix and its column-vector representation:
\[
  \bm{A}
  = \begin{bmatrix}
      a_{00} & a_{01} & \cdots & a_{0,n-1} \\
      a_{10} & a_{11} & \cdots & a_{1,n-1} \\
      \vdots & \vdots & \ddots & \vdots    \\
      a_{n-1,0} & \cdots & \cdots & a_{n-1,n-1}
    \end{bmatrix}
  = \begin{bmatrix}
      \bm{a}_0 & \bm{a}_1 & \cdots & \bm{a}_{n-1}
    \end{bmatrix}.
\]

The fundamental algebraic operations are:
\begin{align*}
  \text{Scalar mult.:}\quad
      &\bm{A}=\gamma\bm{B} \;\Rightarrow\; a_{ij} = \gamma b_{ij}.\\[2pt]
  \text{Addition:}\quad
      &\bm{A}=\bm{B}\pm\bm{C} \;\Rightarrow\; a_{ij} = b_{ij}\pm c_{ij}.\\[2pt]
  \text{Matrix-vector product:}\quad
      &\bm{y}=\bm{A}\bm{x} \;\Rightarrow\;
          y_i = \sum_{j=0}^{n-1} a_{ij}\,x_j.\\[2pt]
  \text{Matrix-matrix product:}\quad
      &\bm{A}=\bm{B}\bm{C} \;\Rightarrow\;
          a_{ij} = \sum_{k=0}^{n-1} b_{ik}\,c_{kj}.\\[2pt]
  \text{Transpose:}\quad
      &\bm{A}=\bm{B}^{T} \;\Rightarrow\; a_{ij} = b_{ji}.
\end{align*}

\paragraph{Inverse.}
If $\bm{A}$ is square and nonsingular, its inverse satisfies
$\bm{A}^{-1}\bm{A}=\bm{I}$, where $\bm{I}$ is the identity matrix.

\paragraph{Hermitian conjugate.}
The Hermitian conjugate $\bm{A}^{\dagger}$ is obtained by transposing
$\bm{A}$ and complex-conjugating every element:
$(A^{\dagger})_{ij} = a_{ji}^{*}$.

\paragraph{Nonsingularity: equivalent characterisations.}
For an $n\times n$ matrix $\bm{A}$ the following properties are all
equivalent: (i)~$\bm{A}^{-1}$ exists; (ii)~$\bm{A}\bm{x}=\bm{0}$
implies $\bm{x}=\bm{0}$; (iii)~the rows of $\bm{A}$ form a basis of
$\mathbb{C}^n$; (iv)~the columns of $\bm{A}$ form a basis of
$\mathbb{C}^n$; (v)~$\bm{A}$ is a product of elementary matrices;
(vi)~zero is not an eigenvalue of $\bm{A}$.

%----------------------------------------------------------------------
\section{Special matrix types}
\label{sec:specialmat}

Table~\ref{tab:matrices} lists the matrix types that appear most
frequently in quantum mechanics and many-body physics.

\begin{table}[h]
\centering
\renewcommand{\arraystretch}{1.35}
\begin{tabular}{lll}
\toprule
\textbf{Name} & \textbf{Defining relation} & \textbf{Element condition} \\
\midrule
Symmetric     & $\bm{A}=\bm{A}^{T}$
              & $a_{ij}=a_{ji}$ \\
Real orthogonal & $\bm{A}=(\bm{A}^{T})^{-1}$
              & $\sum_k a_{ik}a_{jk}=\delta_{ij}$ \\
Hermitian     & $\bm{A}=\bm{A}^{\dagger}$
              & $a_{ij}=a_{ji}^{*}$ \\
Unitary       & $\bm{A}=(\bm{A}^{\dagger})^{-1}$
              & $\sum_k a_{ik}a_{jk}^{*}=\delta_{ij}$ \\
Real matrix   & $\bm{A}=\bm{A}^{*}$
              & $a_{ij}=a_{ij}^{*}$ \\
\bottomrule
\end{tabular}
\caption{Important special matrix types and their defining properties.}
\label{tab:matrices}
\end{table}

Structured (sparse) matrices important for numerical many-body work:
\begin{itemize}
  \item \textbf{Diagonal}: $a_{ij}=0$ for $i\neq j$.
  \item \textbf{Upper/lower triangular}: $a_{ij}=0$ for $i>j$ / $i<j$.
  \item \textbf{Upper/lower Hessenberg}: $a_{ij}=0$ for $i>j+1$ / $i<j+1$.
  \item \textbf{Tridiagonal}: $a_{ij}=0$ for $|i-j|>1$.
  \item \textbf{Banded with bandwidth $p$}: $a_{ij}=0$ for $|i-j|>p$.
  \item Block upper/lower triangular matrices, and more.
\end{itemize}

\begin{notebox}
\textbf{Many-body connection.} Hamiltonian matrices in shell-model
(full configuration-interaction, FCI) calculations are typically large,
sparse, and often banded.  Exploiting sparsity is essential to
diagonalise them for systems with many single-particle states.
\end{notebox}

%----------------------------------------------------------------------
\section{Dirac bra-ket notation}
\label{sec:braket}

In quantum mechanics it is conventional to write state vectors using
Dirac's bra-ket formalism.  A ket $\bket{x}$ is a column vector
\[
  \bket{x} =
  \begin{bmatrix} x_0 \\ x_1 \\ \vdots \\ x_{n-1} \end{bmatrix},
\]
and the corresponding bra $\bbra{x}$ is the row vector (dual)
\[
  \bbra{x} =
  \begin{bmatrix} x_0^{*} & x_1^{*} & \cdots & x_{n-1}^{*} \end{bmatrix},
\]
so $\bket{x}^{\dagger}=\bbra{x}$.  The \emph{inner product}
(a complex number) is
\[
  \braket{y}{x} = \sum_{i=0}^{n-1} y_i^{*} x_i.
\]
Note that $\braket{x}{x}\geq 0$ is always real and non-negative,
and $\braket{y}{x}^{*}=\braket{x}{y}$.

The \emph{outer product} $\ketbra{x}{y}$ is an $n\times n$ matrix
with elements $(\ketbra{x}{y})_{ij}=x_i y_j^{*}$:
\[
  \bket{x}\bbra{y} =
  \begin{bmatrix}
    x_0 y_0^{*} & x_0 y_1^{*} & \cdots & x_0 y_{n-1}^{*} \\
    x_1 y_0^{*} & x_1 y_1^{*} & \cdots & x_1 y_{n-1}^{*} \\
    \vdots      & \vdots      & \ddots & \vdots           \\
    x_{n-1}y_0^{*} & \cdots & \cdots & x_{n-1}y_{n-1}^{*}
  \end{bmatrix}.
\]

\paragraph{Worked example.}
Let $\bket{x}=\begin{bmatrix}1-\imath \\ 2+\imath\end{bmatrix}$.
Then
\[
  \braket{x}{x}
  = (1+\imath)(1-\imath)+(2-\imath)(2+\imath)=2+5=7\in\mathbb{R}.
\]
For two different vectors
$\bket{x}=\begin{bmatrix}-1\\2\imath\\1\end{bmatrix}$,
$\bket{y}=\begin{bmatrix}1\\0\\\imath\end{bmatrix}$:
$\braket{x}{y}=-1+\imath\neq\braket{y}{x}=-1-\imath$,
confirming the rule $\braket{x}{y}^{*}=\braket{y}{x}$.

%----------------------------------------------------------------------
\section{Computational basis states and superposition}
\label{sec:compbasis}

A \emph{two-level system} (qubit) is described by two orthonormal
basis states, which we label
\[
  \bket{0}=\begin{bmatrix}1\\0\end{bmatrix},
  \qquad
  \bket{1}=\begin{bmatrix}0\\1\end{bmatrix},
  \qquad
  \braket{0}{1}=0,\quad
  \braket{0}{0}=\braket{1}{1}=1.
\]
These are the \emph{computational basis} states.  A general state is
a \emph{superposition}
\[
  \bket{\psi} = \alpha\bket{0}+\beta\bket{1},
\]
where $\alpha=\braket{0}{\psi}$ and $\beta=\braket{1}{\psi}$ are
the probability amplitudes, satisfying
$|\alpha|^{2}+|\beta|^{2}=1$.

The identity operator decomposes as a sum of outer products:
\[
  \bm{I}
  = \sum_{i=0}^{1}\ketbra{i}{i}
  = \begin{bmatrix}1&0\\0&0\end{bmatrix}
   +\begin{bmatrix}0&0\\0&1\end{bmatrix}
  = \begin{bmatrix}1&0\\0&1\end{bmatrix}.
\]

\paragraph{Projection operators.}
The outer products
\[
  \bm{P}=\ketbra{0}{0}=\begin{bmatrix}1&0\\0&0\end{bmatrix},
  \qquad
  \bm{Q}=\ketbra{1}{1}=\begin{bmatrix}0&0\\0&1\end{bmatrix}
\]
are \emph{projection operators}.  They are \emph{idempotent}
($\bm{P}^2=\bm{P}$, $\bm{Q}^2=\bm{Q}$), commute with each other
($[\bm{P},\bm{Q}]=0$), and satisfy $\bm{P}\bm{Q}=0$,
$\bm{P}+\bm{Q}=\bm{I}$.  Acting on a superposition state:
\[
  \bm{P}\bket{\psi}=\alpha\bket{0},\qquad
  \bm{Q}\bket{\psi}=\beta\bket{1}.
\]
The measurement postulate states that a measurement on $\bket{\psi}$
yields $\bket{0}$ with probability $|\alpha|^2$ and $\bket{1}$ with
probability $|\beta|^2$.

\begin{notebox}
\textbf{Many-body connection.} In many-body theory, projection
operators $\hat{P}$ (onto occupied states) and $\hat{Q}$ (onto
unoccupied / virtual states) define the particle-hole decomposition
that is fundamental to perturbation theory, Hartree-Fock theory, and
coupled-cluster methods.
\end{notebox}

\paragraph{Density matrix.}
The density matrix (density operator) of a pure state $\bket{\psi}$
is the outer product
\[
  \bm{\rho} = \ketbra{\psi}{\psi}
  = \begin{bmatrix}\alpha\alpha^{*}&\alpha\beta^{*}\\\beta\alpha^{*}&\beta\beta^{*}\end{bmatrix},
\]
with $\mathrm{tr}\,\bm{\rho}=|\alpha|^2+|\beta|^2=1$.  Density matrices
generalise to mixed states and subsystems; they are central to quantum
information theory and to reduced-density-matrix many-body methods.

%----------------------------------------------------------------------
\section{Pauli matrices}
\label{sec:pauli}

The three Pauli matrices are
\[
  \sigma_x = \begin{bmatrix}0&1\\1&0\end{bmatrix},\quad
  \sigma_y = \begin{bmatrix}0&-\imath\\\imath&0\end{bmatrix},\quad
  \sigma_z = \begin{bmatrix}1&0\\0&-1\end{bmatrix}.
\]
Key properties (easy to verify directly):
\begin{itemize}
  \item \textbf{Involutory}: $\sigma_x^2=\sigma_y^2=\sigma_z^2=\bm{I}$.
  \item \textbf{Unitary}: $\sigma_\alpha^\dagger=\sigma_\alpha^{-1}$
        (Hermitian conjugate equals inverse).
  \item \textbf{Hermitian}: $\sigma_\alpha^\dagger=\sigma_\alpha$.
  \item $\det(\sigma_\alpha)=-1$ for each $\alpha$.
  \item \textbf{Commutation}: $[\sigma_x,\sigma_y]=2\imath\sigma_z$
        (and cyclic permutations).
\end{itemize}

The Pauli matrices act on the computational basis as
\[
  \sigma_x\bket{0}=\bket{1},\quad
  \sigma_x\bket{1}=\bket{0};
  \qquad
  \sigma_z\bket{0}=+\bket{0},\quad
  \sigma_z\bket{1}=-\bket{1}.
\]
Acting on a superposition $\bket{\psi}=\alpha\bket{0}+\beta\bket{1}$:
\[
  \sigma_x\bket{\psi}=\beta\bket{0}+\alpha\bket{1},\qquad
  \sigma_y\bket{\psi}=\imath\beta\bket{0}-\imath\alpha\bket{1}.
\]

\begin{notebox}
\textbf{Quantum computing connection.} The Pauli matrices are the
elementary single-qubit gates $X$, $Y$, $Z$.  Together with the
Hadamard gate $H=(\sigma_x+\sigma_z)/\sqrt{2}$ and the phase gate $T$
they form a universal single-qubit gate set.

\textbf{Many-body connection.} Via the Jordan-Wigner transformation,
the creation/annihilation operators of second quantisation are mapped
directly onto products (strings) of Pauli matrices, enabling the
representation of fermionic Hamiltonians on qubit registers.
\end{notebox}

%----------------------------------------------------------------------
\section{Orthonormal bases and the completeness relation}
\label{sec:onb}

A set of states $\{\bket{\phi_i}\}_{i=0}^{n-1}$ is an
\emph{orthonormal basis} (ONB) of an $n$-dimensional Hilbert space if
\[
  \braket{\phi_j}{\phi_i}=\delta_{ij}
  \quad\text{and}\quad
  \sum_{i=0}^{n-1}\ketbra{\phi_i}{\phi_i}=\bm{I}
  \quad\text{(completeness / resolution of the identity)}.
\]
Any state can then be expanded as
$\bket{\psi}=\sum_i\braket{\phi_i}{\psi}\bket{\phi_i}$,
with $\braket{\phi_i}{\psi}$ the projection (overlap) of $\bket{\psi}$
onto $\bket{\phi_i}$.

\begin{notebox}
\textbf{Many-body connection.} The natural ONB in many-body physics
is the set of single-particle eigenstates of the one-body Hamiltonian:
\[
  \hat{h}_0\bket{\phi_i}=\varepsilon_i\bket{\phi_i},
  \qquad
  \phi_i(\mathbf{r})=\braket{\mathbf{r}}{\phi_i}.
\]
This single-particle basis $\{\bket{\phi_i}\}$ is the starting point
for virtually all many-body methods: Hartree-Fock, MBPT, FCI,
coupled-cluster, and quantum Monte Carlo.
\end{notebox}

%----------------------------------------------------------------------
\section{Unitary transformations and preservation of orthogonality}
\label{sec:unitary}

An $n\times n$ matrix $\bm{U}$ is \emph{unitary} if
$\bm{U}^{\dagger}\bm{U}=\bm{U}\bm{U}^{\dagger}=\bm{I}$, which means
$\bm{U}^{-1}=\bm{U}^{\dagger}$.

\paragraph{Orthogonality is preserved.}
If $\{\bm{v}_i\}$ is an orthonormal set, $\bm{v}_j^T\bm{v}_i=\delta_{ij}$,
then the transformed set $\bm{w}_i=\bm{U}\bm{v}_i$ is also orthonormal:
\[
  \bm{w}_j^T\bm{w}_i
  =(\bm{U}\bm{v}_j)^T\bm{U}\bm{v}_i
  =\bm{v}_j^T\underbrace{\bm{U}^T\bm{U}}_{=\,\bm{I}}\bm{v}_i
  =\bm{v}_j^T\bm{v}_i=\delta_{ij}.
\]
In Dirac notation, if $\bket{\psi_i}=\sum_j u_{ij}\bket{\phi_j}$
with $(u_{ij})$ unitary, then $\braket{\psi_j}{\psi_i}=\delta_{ij}$.

\paragraph{Why only unitary transformations?}
For the inner product (and hence the norm) to be preserved under
$\bket{\phi}=\bm{U}\bket{\psi}$ we need
\[
  \braket{\phi}{\phi}
  =\bbra{\psi}\bm{U}^{\dagger}\bm{U}\bket{\psi}
  =\braket{\psi}{\psi}=1,
\]
which requires $\bm{U}^{\dagger}\bm{U}=\bm{I}$.  Since probabilities
must sum to unity, this is a fundamental physical requirement.
Unitary transformations are rotations (and reflections) in Hilbert
space.  The fact that $\bm{U}^{-1}=\bm{U}^{\dagger}$ also means that
every unitary transformation is perfectly reversible.

\paragraph{Example: two-state rotation.}
Given $\bket{\psi}=\alpha\bket{0}+\beta\bket{1}$ and a unitary $\bm{U}$,
the rotated state
$\bket{\phi}=\gamma\bket{0}+\delta\bket{1}=\bm{U}\bket{\psi}$ satisfies
\[
  \begin{bmatrix}\gamma\\\delta\end{bmatrix}
  =\begin{bmatrix}u_{00}&u_{01}\\u_{10}&u_{11}\end{bmatrix}
  \begin{bmatrix}\alpha\\\beta\end{bmatrix}.
\]
Since $|\alpha|^2+|\beta|^2=1$ and $\bm{U}$ is unitary,
$|\gamma|^2+|\delta|^2=1$ automatically.

\begin{notebox}
\textbf{Many-body connection.} In Hartree-Fock theory the
self-consistent field (SCF) procedure generates a sequence of unitary
rotations of the single-particle basis.  The total HF energy is
invariant under unitary rotations within the occupied or virtual
subspace separately, which is the Thouless theorem.

\textbf{Quantum computing connection.} Every quantum gate is a
unitary matrix.  The evolution of a closed quantum system is described
entirely by unitary operations; this is the quantum analogue of
classical reversibility.
\end{notebox}

%----------------------------------------------------------------------
\section{Tensor products}
\label{sec:tensor}

The \emph{tensor (Kronecker) product} of two vectors $\bket{x}$
(length $n$) and $\bket{y}$ (length $m$) is the vector of length $nm$:
\[
  \bket{x}\otimes\bket{y}
  =\bket{xy}
  =\begin{bmatrix}x_0 y_0\\ x_0 y_1\\ x_1 y_0\\ x_1 y_1\end{bmatrix}
  \quad\text{(for } n=m=2\text{)}.
\]
For the computational basis:
\[
  \bket{0}\otimes\bket{0}=\bket{00}
    =\begin{bmatrix}1\\0\\0\\0\end{bmatrix},\quad
  \bket{0}\otimes\bket{1}=\bket{01}
    =\begin{bmatrix}0\\1\\0\\0\end{bmatrix},\quad
  \bket{1}\otimes\bket{0}=\bket{10}
    =\begin{bmatrix}0\\0\\1\\0\end{bmatrix},\quad
  \bket{1}\otimes\bket{1}=\bket{11}
    =\begin{bmatrix}0\\0\\0\\1\end{bmatrix}.
\]
These four states form an ONB for the two-qubit Hilbert space
$\mathbb{C}^4=\mathbb{C}^2\otimes\mathbb{C}^2$.

The most general state of $N$ qubits is
\[
  \bket{\Psi}
  = \sum_{i_1 i_2\cdots i_N}\Psi_{i_1 i_2\cdots i_N}
    \bket{i_1}\otimes\bket{i_2}\otimes\cdots\otimes\bket{i_N},
\]
requiring $2^N$ complex amplitudes.  For $N\approx 300$ this
exceeds the number of atoms in the observable universe — this is the
\emph{exponential complexity} at the heart of quantum many-body physics.

\begin{notebox}
\textbf{Many-body connection.} The full $N$-particle Hilbert space
is a tensor product of $N$ single-particle spaces.  The FCI method
works in this exponentially large space; all approximate methods
(HF, MBPT, coupled cluster, DMRG, \ldots) are strategies to identify
the relevant low-dimensional submanifold.

\textbf{Quantum computing connection.} A quantum computer with $N$
qubits \emph{inherently} represents a $2^N$-dimensional state vector.
This is both the source of its potential power and the reason that
quantum simulation of many-body systems is a leading application.
\end{notebox}

%----------------------------------------------------------------------
\section{Spectral decomposition of Hermitian operators}
\label{sec:spectral}

A Hermitian operator $\bm{A}=\bm{A}^{\dagger}$ has the following key
properties:
\begin{enumerate}
  \item All eigenvalues $\lambda_i$ are \emph{real}.
  \item Eigenvectors belonging to distinct eigenvalues are
        \emph{orthogonal}.
  \item $\bm{A}$ is diagonalisable (it is a \emph{normal} operator,
        $[\bm{A},\bm{A}^{\dagger}]=0$).
\end{enumerate}

The eigenvalue equation $\bm{A}\bket{\psi}=\lambda\bket{\psi}$
has non-trivial solutions only when
\[
  \det(\bm{A}-\lambda\bm{I})=0
  \quad\text{(characteristic polynomial).}
\]

\paragraph{Spectral decomposition.}
Let $\{\bket{i}\}$ be the eigenbasis of $\bm{A}$ with eigenvalues
$\lambda_i$.  Any state expands as
$\bket{\psi}=\sum_i\alpha_i\bket{i}$.
Using projection operators $\bm{P}_i=\ketbra{i}{i}$ we obtain:
\[
  \bm{A}\bket{\psi}
  =\sum_i\alpha_i\lambda_i\bket{i}
  =\sum_i\lambda_i\bm{P}_i\bket{\psi},
\]
from which, since this holds for any $\bket{\psi}$,
\[
  \boxed{\bm{A}=\sum_{i=0}^{n-1}\lambda_i\,\ketbra{i}{i}
               =\sum_{i=0}^{n-1}\lambda_i\,\bm{P}_i.}
\]
This is the \emph{spectral decomposition}: the operator is expressed
as a weighted sum of its projection operators.

\paragraph{Example for a $2\times 2$ operator.}
If $\bm{A}$ has eigenstates
$\bket{\psi_a}=\alpha_0\bket{0}+\alpha_1\bket{1}$ and
$\bket{\psi_b}=\beta_0\bket{0}+\beta_1\bket{1}$
with eigenvalues $\lambda_a$, $\lambda_b$, then
\[
  \bm{A}
  =\lambda_a\begin{bmatrix}|\alpha_0|^2 & \alpha_0\alpha_1^*\\
                            \alpha_1\alpha_0^* & |\alpha_1|^2\end{bmatrix}
  +\lambda_b\begin{bmatrix}|\beta_0|^2 & \beta_0\beta_1^*\\
                            \beta_1\beta_0^* & |\beta_1|^2\end{bmatrix}.
\]

\begin{notebox}
\textbf{Many-body connection.} The spectral decomposition is used
every time one diagonalises a Hamiltonian matrix.  In FCI the full
$N$-body Hamiltonian is diagonalised in the many-determinant basis;
in Hartree-Fock the Fock matrix (a single-particle Hermitian operator)
is diagonalised to find the optimal orbital basis.

\textbf{Quantum computing connection.} Any unitary gate $\bm{U}$ can
be written as $\bm{U}=e^{i\bm{H}}$ for a Hermitian generator $\bm{H}$.
The spectral decomposition of $\bm{H}$ gives
$\bm{U}=\sum_i e^{i\lambda_i}\ketbra{i}{i}$.
This exponential relationship connects the physical Hamiltonian to the
time-evolution operator $\hat{U}=e^{-i\hat{H}t/\hbar}$ and is the
basis of the variational quantum eigensolver (VQE) and unitary
coupled-cluster (UCC) ansatz.
\end{notebox}

